%% kapa: exam=1 round=19 year=2008 period=1 form=A subject=accounting subject_ko=회계학 seq=17 number=97
%% source: Q-Net 감정평가사 기출문제 (cst00309) artlSeq=1001962 file=extracted/1008/1001962/2008년 제19회 감정평가사 1차시험 기출문제/1교시.hwp
%% conversion: kordoc 4.24.51 -> tools/gen1.py (auto); review=proofread-sonnet; flags=none
\begin{question}{97}
유가증권에 대한 다음의 설명 중 가장 옳은 것은?
\begin{choices}
  \item 시장성을 상실하여 단기매매증권에서 매도가능증권으로 재분류된 유가증권은 시장성을 회복할 경우 다시 단기매매증권으로 재분류할 수 있다.
  \item 단기매매증권의 장부가액과 공정가액의 차이는 당기손익에 반영하지만 매도가능증권의 장부가액과 공정가액의 차이는 자본조정으로 처리한다.
  \item 매도가능증권손상차손환입은 기타포괄손익누계액 항목으로 보고되며 손익계산서에 반영되지 않는다.
  \item 투자채권의 경우 매도가능증권에서 만기보유증권으로 재분류는 허용되지만 만기보유증권에서 매도가능증권으로의 재분류는 인정되지 않는다.
  \item 투자채권 중 만기보유증권은 기준금리의 상승에 따라 가격이 하락하는 경우에도 공정가치변동을 인식하지 않는다.
\end{choices}
\end{question}
\answer{5}
